\section{Introduction}

A typed decision model is given a state and a question with declared answers,
and returns a probability for each answer instead of generating text
\citep{pai2026jevbench}. Models of this kind are sold as guards: they decide
whether a prompt is a jailbreak, whether a text tries to override a system's
instructions, whether a chatbot's reply is harmful, and whether an agent acted
unsafely. A guard sits in front of other systems, so its decision is an action:
a request passes or is blocked.

The users of these systems do not all write in English, and neither do
attackers. Work on generative models has shown that a harmful request
translated into a low-resource language can pass safety training that holds in
English \citep{yong2023lowresource,deng2023multilingualjailbreak,wang2023alllanguages}.
Whether the same gap exists for decision models has not been measured. Their
public evaluations, including our own JevBench \citep{pai2026jevbench}, use
English data only.

This paper measures it. We translate the five safety slices of JevBench into
eight languages and run \MLnModels{} hosted decision models on every language and
on the English originals. We follow the rule that a result names its language
\citep{ducel2022benderrule}: every number below is for a named language, and
English is one of them.

We make four contributions.

\begin{enumerate}
  \item \textbf{A translated safety benchmark.} \MLnCasesEN{} English cases from
  five public safety datasets, translated into Chinese, Arabic, Hindi,
  Vietnamese, Swahili, Spanish, French and Russian, covering five scripts and
  five language families (Section~\ref{sec:benchmark}). Every
  translated string passes automatic checks, and a case with a failed string is
  left out rather than repaired.
  \item \textbf{A paired comparison with English.} Each model is compared with
  its own English answers on exactly the cases its translated run answered, with
  bootstrap intervals over case families, so a language is never credited or
  penalised for the cases its translation dropped (Section~\ref{sec:setup}).
  \item \textbf{A separation of two failure modes.} By reporting AUROC and the
  share of positive answers next to accuracy, we show where a model loses the
  ability to tell harmful from harmless content and where it keeps that ability
  but moves its probabilities across the decision threshold
  (Section~\ref{sec:threshold}).
  \item \textbf{A case-level audit.} We list the cases with the largest change
  between English and each language, and find both genuine model failures and
  changes made by the translation itself (Section~\ref{sec:examples}).
\end{enumerate}

Our aim is diagnostic, as in JevBench. We do not rank the models for production
use, and Section~\ref{sec:limitations} sets out why. The results and every
example are also browsable at
\url{https://leanmcp.github.io/jevbench/multilingual.html}.
